% エントリ 10(2026-11-07): 方針の見直し
\section{[E10, 2026-11-07] 方針の見直し: GW のスピンだけで閉じた解析}\label{sec:e10}

E9 のあと、解析の方針を四点見直した。計算は \texttt{calc/e10\_gwonly.py}、\texttt{calc/hbr16.py}(合体残骸のスピン)、\texttt{calc/chisr.py}(集団解析用の $\chi_\mathrm{SR}$)、出力は \texttt{calc/out/e10.json}。

\subsection{主解析は GW のスピンだけで閉じる}

X 線連星のスピンは、スペクトルの模型による系統誤差が大きい(E9)。
\begin{itemize}
\item Cyg X-1 は、同じデータから $>0.9985$ にも $0.87$--$0.92$ にもなり、暖かい Compton 層を入れると $\lesssim0.1$ になる。
\item GRS 1915+105 は、統計誤差だけで $0.98\pm0.01$、系統誤差を入れた 3$\sigma$ では $>0.92$ である。
\end{itemize}
そこで主解析は「自己相互作用を入れて、LVK のスピンの測定を自己無撞着に解析し直す」形にする。X 線連星は付録で整合性の確認として扱う。これに伴い、窓の上限は GW241011 の主星だけで決まる。

\begin{figure}[H]
\centering
% E10: (a) GW241011 だけで決めた窓で回転を失える 1 世代目の最小質量、(b) 2 世代目の生まれたときのスピン。データは calc/e10_gwonly.py
\begin{tikzpicture}
\begin{groupplot}[group style={group size=2 by 1, horizontal sep=18mm}, width=0.46\linewidth, height=0.42\linewidth, note]
\nextgroupplot[xmode=log, xmin=0.45e-12, xmax=1.1e-12, ymin=6, ymax=23,
  xtick={0.5e-12,0.6e-12,0.7e-12,0.8e-12,1e-12}, xticklabels={0.5,0.6,0.7,0.8,1.0},
  xlabel={$\mu\ [10^{-12}\,\mathrm{eV}]$}, ylabel={$M_\mathrm{min}\ [M_\odot]$},
  title={\small (a) 回転を失える 1 世代目の最小質量},
  legend style={font=\scriptsize, at={(0.98,0.98)}, anchor=north east, fill=white, fill opacity=0.85, text opacity=1, draw=none}, legend cell align=left]
  \fill[cB!10] (axis cs:0.60e-12,6) rectangle (axis cs:0.89e-12,23);
  \addplot[cA!50, line width=0.9pt] table[x=mu, y=opt_d18] {figs/data/e10_mmin.dat}; \addlegendentry{楽観的、$\delta=0.18$}
  \addplot[cA!50, line width=0.9pt, dashed] table[x=mu, y=opt_d05] {figs/data/e10_mmin.dat}; \addlegendentry{楽観的、$\delta=0.05$}
  \addplot[cA, line width=0.9pt] table[x=mu, y=real_d18] {figs/data/e10_mmin.dat}; \addlegendentry{現実的、$\delta=0.18$}
  \addplot[cA, line width=0.9pt, dashed] table[x=mu, y=real_d05] {figs/data/e10_mmin.dat}; \addlegendentry{現実的、$\delta=0.05$}
  \addplot[ink, line width=0.8pt, densely dotted] table[x=mu, y=realGRS_d18] {figs/data/e10_mmin.dat}; \addlegendentry{現実的 + GRS、$\delta=0.18$}
\nextgroupplot[xmin=0.3, xmax=0.95, ymin=0, ymax=22,
  xlabel={生まれたときのスピン $\chi$}, ylabel={確率密度},
  title={\small (b) 2 世代目の生まれたときのスピン},
  legend style={font=\scriptsize, at={(0.02,0.98)}, anchor=north west, fill=white, fill opacity=0.85, text opacity=1, draw=none}, legend cell align=left]
  \addplot[cM, line width=0.9pt, const plot mark mid] table[x=chi, y=p0] {figs/data/e10_natal.dat}; \addlegendentry{親が回らない}
  \addplot[cC, line width=0.9pt, const plot mark mid] table[x=chi, y=p1] {figs/data/e10_natal.dat}; \addlegendentry{$\mathcal U(0,0.5)$}
  \addplot[cD, line width=0.9pt, const plot mark mid] table[x=chi, y=p2] {figs/data/e10_natal.dat}; \addlegendentry{$\mathrm{Beta}(2,5)$}
  \addplot[cD!50, line width=0.9pt, dashed, const plot mark mid] table[x=chi, y=p3] {figs/data/e10_natal.dat}; \addlegendentry{$\mathrm{Beta}(2,5)$、$q\ge0.25$}
  \draw[ink2, densely dotted] (axis cs:0.69,0) -- (axis cs:0.69,22);
  \draw[cA, densely dotted] (axis cs:0.72,0) -- (axis cs:0.72,22);
  \draw[cB, densely dotted] (axis cs:0.78,0) -- (axis cs:0.78,22);
  \node[font=\scriptsize, color=cB, anchor=south west, rotate=90] at (axis cs:0.785,8) {0.78(観測値)};
  \node[font=\scriptsize, color=cA, anchor=south west, rotate=90] at (axis cs:0.725,8) {0.72};
\end{groupplot}
\end{tikzpicture}

\caption{(a) GW241011 の主星だけで $f$ の上限を決めたとき、6 Gyr のうちに 211 の臨界 $+\delta$ まで回転を失える 1 世代目の最小質量($\chi_0=0.7$)。点線は E9 の現実的な GRS 1915+105 の上限も入れた場合。橙の帯は \arxiv{2609.02678} の質量(90\%)。(b) 2 世代目の生まれたときのスピン(\arxiv{1605.01938} の一般スピン版)。親の質量比は $q\in[0.5,1]$ で一様(破線は $[0.25,1]$)、親のスピンの向きは等方。}
\label{fig:e10}
\end{figure}

\begin{result}
$\mu=0.71\times10^{-12}\eV$ では次のようになる(主要項の $\gamma_\mathrm{BH}$)。
\begin{itemize}
\item \textbf{GW241011 だけの上限}: 楽観的な仮定(0.78 で生まれ、$T_2=3\times10^7\yr$)で $f\le2.4\times10^{12}\GeV$、現実的な仮定(0.72 で生まれ、$5\times10^7\yr$)で $1.1\times10^{12}\GeV$。
\item \textbf{回転を失える 1 世代目の最小質量}($\delta=0.18/0.05$):
  \begin{itemize}
  \item 楽観的: $8.8/10.9\Msun$
  \item 現実的: $11.7/14.5\Msun$
  \item 現実的な仮定に GRS 1915+105 を加えると: $15.5/19.3\Msun$
  \end{itemize}
  $\mu=0.60$、$0.89\times10^{-12}\eV$ でも 1--2 $\Msun$ しか動かない(図~\ref{fig:e10}(a))。
\item \textbf{判定の対象になる質量の上限}: 211 の臨界 $+\delta$ が初期スピン 0.7 に届く質量は、$\delta=0.18/0.05$ で $26/35\Msun$($\mu=0.71\times10^{-12}\eV$)。これより重い 1 世代目は、この判定では扱えない。
\item GW のスピンだけなら、現実的な仮定でも $\sim12$--$15\Msun$ 以上の 1 世代目は回転を失える。GRS 1915+105 を現実的な年齢で入れても $15$--$19\Msun$ 以上である。E9 の「$\sim20\Msun$ 以上だけ」は、1 世代目を 15・20・25 $\Msun$ の格子でしか見ていなかったための粗い言い方だった(訂正)。
\end{itemize}
\end{result}

\subsection{図の曲線の定義}

E8・E9 と論文の図の曲線は、次のように定義した。
\begin{itemize}
\item \textbf{上限の曲線}(GW241011、GRS 1915+105): BH の(生まれたときのスピン、許すスピンの減少、年齢)を与え、その減少が年齢より長くかかる $f$ の上限を $\mu$ の関数として描いた。$f$ が大きいほど回転を失う速さが $f^2$ で速くなるので、上限になる。
\item \textbf{下限の曲線}(1 世代目): 質量を与え、生まれたときのスピン 0.7 から「211 の臨界 $+\delta$」まで、年齢 6 Gyr のうちに落ちる $f$ の下限を描いた。
  \begin{itemize}
  \item これより下では自己相互作用が強すぎて、その質量の 1 世代目が回転を失わない。そのため、質量とスピンの相関を超放射で説明できない。
  \item $\delta=0.18$ は \arxiv{2609.02678} の積み上がりの幅 $\sigma_\mathrm{SR}$ にあたる。
  \item 6 Gyr は、球状星団の低赤方偏移の 1 世代目の年齢の目安(E8、ただし E9 の注意つき)である。
  \end{itemize}
\end{itemize}
どちらも集団解析ではなく説明のための判定で、本解析では年齢・自然スピンを分布として周辺化する。

\subsection{2 世代目の生まれたときのスピン}

E9 の「現実的な 0.72」は、球状星団の模擬計算の 90\% 区間の上端(\arxiv{1906.10260} の Fig.~8、1 世代目のスピン 0.5 の場合)から取った値だった。合体残骸のスピンの当てはめ(\arxiv{1605.01938} の式 13--17、Table 1 の上段)を実装し、自分で分布を出した。
\begin{itemize}
\item \textbf{当てはめの照合}: 等質量・無スピンで 0.68646(式 11 の拘束)、$q=1/3$ で 0.540、等質量で 0.5 ずつ揃ったスピンで 0.832、質量比 $\to0$ で親のスピンに戻る。$k_{00}=-3.821$(原文は $\approx-3.82$)。
\item \textbf{分布}(図~\ref{fig:e10}(b)):
  \begin{itemize}
  \item 親が回らないとき: 中央値 0.675、90\% が 0.63--0.69。
  \item 親のスピンが $\mathcal U(0,0.5)$・$\mathrm{Beta}(2,5)$ で向きが等方のとき: 中央値 0.676--0.677、95 パーセンタイル 0.75--0.76。
  \item 0.72 以上になる確率は 15--18\%、0.78 以上は 0.5--1.8\%。
  \item 質量比を $q\ge0.25$ に広げると、中央値は 0.657 に下がる。
  \end{itemize}
\item \textbf{意味}:
  \begin{itemize}
  \item 0.72 は分布の上位 15--18\% の値で、「現実的」としてはやや甘い。
  \item 0.78(観測値)で生まれる確率は 2\% 以下である。
  \item 典型的な 2 世代目(0.68)は、GW241011 の主星の観測の 90\% 区間(0.69--0.87)の下端のすぐ外にある。観測値が高めに出ているか、親が速く回っていたかのどちらかである。
  \item 本解析では、固定した値ではなく、親の分布から出したこの自然スピンの分布と、観測の事後分布を組み合わせて判定する。
  \end{itemize}
\end{itemize}

\subsection{集団解析用の $\chi_\mathrm{SR}$}

集団解析では、モンテカルロの標本ごとに $\chi_\mathrm{SR}(M,\chi_0,T,\mu,f)$ を計算するので、速くなければならない。準平衡では(E3)
\begin{equation}
\frac{\mathrm{d}\chi}{\mathrm{d}t}=\mu\Bigl(\frac{f}{\Mpl}\Bigr)^2H(\alpha,\chi),\qquad
H=-\frac{2}{\sqrt3}\frac{g^{3/2}\kappa_\infty^{1/2}}{\kappa_\mathrm{BH}\,\alpha^7\,\tilde r_+},\qquad g\equiv\gamma_\mathrm{SR}/\mu
\label{eq:e10_H}
\end{equation}
となり、$f$ と $\mu$ は前の因子にしか入らない。そこで $I(\alpha,\chi)=\int_\chi^{0.995}\mathrm{d}\chi'/|H|$ を一度だけ表にすれば、年齢 $T$ のあとのスピンは
\[
I(\alpha,\chi_T)=I(\alpha,\chi_0)+\tau,\qquad \tau\equiv\mu(f/\Mpl)^2T
\]
を $\chi_T$ について解くだけで決まる。$\chi_\mathrm{SR}$ は $(\alpha,\chi_0,\tau)$ の三つの量だけの関数である。

\begin{result}
\begin{itemize}
\item \textbf{表}: $\alpha\in[0.02,0.30]$ の 141 点と、臨界からの距離 $u$ の 460 点。$\ln I$ と、その逆引きの表を持つ。作るのに 0.9 秒。
\item \textbf{照合}: 無作為な 252 点で、表から引いたスピンと直接の積分(臨界からの距離の対数で変数変換)の差は、中央値 $5.5\times10^{-5}$、最大 $1.8\times10^{-4}$。GW241011 の主星($f=4\times10^{12}\GeV$)で、$T=10^7$、$10^8\yr$ のスピンは 0.6950、0.4490。E5 の ODE の 0.6951、0.4489 と合う。
\item \textbf{計算時間}: $10^6$ 回の評価が 0.21 秒(numpy)。集団解析の中で毎回呼んでも重くない。
\item \textbf{有効な範囲}: 準平衡が速く成り立つ $f$(E5 で $f\lesssim10^{15}\GeV$)と二準位の $\alpha\lesssim0.15$--$0.2$。自由場の側($f\gtrsim f_\mathrm{BC}$)は E2 の半解析(準位の飽和で区切る方法)を使い、両者をつなぐ扱いは本解析の前に決める。
\end{itemize}
\end{result}

\begin{caveat}
\begin{itemize}
\item 自己相互作用の率は主要項で、相対論的な補正(E7)と BH の質量の変化(E9)は表に入れていない。どちらも \texttt{chisr.H} の引数で入れられる。
\item 2 世代目の自然スピンの分布は、親の質量比とスピンの分布の仮定による。本解析では親の分布を 1 世代目の集団の推定と同時に決める。
\end{itemize}
\end{caveat}
